\chapter{THE RATIONALS}

\textit{July 20th.}

Prerequisites for this course include Rings and Modules (tensor products, localization, etc.), Field and Galois Theory, and bits of Commutative Algebra and Algebraic Geometry.

\section{Algebraic Extensions}
The basic object of study in algebraic number theory is the field of rational numbers $\Q$. We recall the field inclusions $\Q \subset \R \subset \C$. We will be interested in the field of algebraic numbers $\overline{\Q} \subset \C$.

\begin{definition}
    We say $k \subseteq L$ is a \eax{field extension} if $k$ and $L$ are both fields. In this case, we say $L$ is an extension of $k$ and write $L/k$. 
\end{definition}

Let $k$ be a field, and $k[X]$ the polynomial ring over $k$. Let $f(X) \in k[X]$ be a polynomial with $\deg(f) = n$, and $\ip{f(x)} \subseteq k[X]$ the ideal generated by $f(X)$. Now define $V \defeq k[X]/\ip{f(X)} = \{g(X) + \ip{f(X)} \mid g(X) \in k[X]\}$. This $V$ is a vector space over $k$. Moreover, $\dim_{k}(V) = n = \deg(f)$; let $\overline{g(X)} = g(\overline{X}) = g(X) + \ip{f(X)}$ denote an element in $V$. Using the divison algorithm, we can write $g = fh + r$ with $\deg(r) < n$, and hence $g(\overline{X}) = r(\overline{X})$. Thus, the set $\cB = \{\overline{1}, \overline{X}, \ldots, \overline{X}^{n-1}\}$ is a generating set for $V$. It is also linearly independent, showing that $\cB$ is a basis for $V$, and hence $\dim_{k}(V) = n$.

\begin{theorem}
    Let $k$ be a field, and $f(X) \in k[X]$ with $\deg f > 0$. There then exists a field $L \supset k$ such that $[L:k] \defeq \dim_{k}L$ is finite, and $f(X)$ has a root in $L$.
\end{theorem}

\begin{proof}
    Without the loss of generality, we may assume that $f(X)$ is irreducible over $k$. Write $f(X) = X^{n} + a_{n-1}X^{n-1} + \cdots + a_{0}$. Consider the companion matrix of $f$:
    \begin{align}
        C(f) = \begin{pmatrix}
            0 & 0 & \cdots & 0 & -a_{0} \\
            1 & 0 & \cdots & 0 & -a_{1} \\
            0 & 1 & \cdots & 0 & -a_{2} \\
            \vdots & \vdots & \ddots & \vdots & \vdots \\
            0 & 0 & \cdots & 1 & -a_{n-1}
        \end{pmatrix}
    \end{align}
    This matrix may be obtained from the $k$-linear map $\ell_{\overline{X}} : k[X]/\ip{f} \to k[X]/\ip{f}$ defined by $\ell_{\overline{X}}(g(\overline{X})) = \overline{X}g(\overline{X})$; the matrix of $\ell_{\overline{X}}$ with respect to the basis $\cB$ formed previously is precisely $C(f)$. Note that the minimal and characteristic polynomials of $C(f)$ are both equal to $f(X)$. Define $\alpha \defeq C(f) \in M_{n}(k)$, and let $L \defeq \{b_{0}+b_{1}\alpha + \cdots + b_{n-1}\alpha^{n-1} \mid b_{i} \in k\} \subseteq M_{n}(k)$. Then $L$ is a ring. In fact, $L$ is a field since any non-zero element's inverse can be computed using the division algorithm. Moreover, $f(\alpha) = 0$ and $L \supset k$.

    As a alternate proof, setting $L \defeq k[X]/\ip{f}$ gives a field extension of $k$ since $\ip{f}$ is a maximal ideal. Moreover, $\overline{X} \in L$ is a root of $f(X)$.
\end{proof}


\begin{enumerate}
    \item Let $R$ be a (commutative) ring, and $\{S_{\lambda}\}_{\lambda \in \Lambda}$ be an arbitrary collection of subrings of $R$. Then $\bigcap_{\lambda \in \Lambda} S_{\lambda}$ is also a subring of $R$.
    \item In particular, the intersection of all subrings of $R$ is a subring of $R$ termed the \eax{prime subring} of $R$ denoted by $\mf{p}_{R}$. It is also the image of the ring homomorphism $\Z \to R$ defined by $n \mapsto n\cdot 1_{R}$.
    \item Let $R \subseteq S$ be rings, and $T \subseteq S$ be a subset. Define $R[T]$ as the intersection of all subrings of $S$ containing $R$ and $T$. It is precisely $\{f(t_{1},\ldots,t_{r}) \mid f \in R[X_{1},\ldots,X_{r}], r \geq 0, t_{i} \in T\}$.
\end{enumerate}
If we let $k \subseteq L$ be fields, and $S \subseteq L$ be a subset, then $k[S]$ is an integral domain. We can then further define $k(S)$ as the field of fractions of $k[S]$, or
\begin{align}
    k(S) = \bigcap_{\substack{E \text{ field} \\ k \subseteq E \subseteq L \\ S \subseteq E}} E = \left\{\frac{f(s_{1},\ldots,s_{r})}{g(s_{1},\ldots,s_{r})} \, \Bigg| \, f,g \in k[X_{1},\ldots,X_{r}], g \neq 0, s_{i} \in S, r \geq 0\right\}.
\end{align}
Thus, $k(S)$ is the smallest subfield of $L$ containing $k$ and $S$. 

\begin{definition}
    An element $\alpha \in \C$ is said to be an \eax{algebraic number} (over $\Q$) if $\alpha$ is a root of a non-zero polynomial over $\Q$
\end{definition}

\begin{definition}
    Define
    \begin{align}
        \overline{\Q} \defeq \{\alpha \in \C \mid \alpha \text{ is an algebraic number}\} \subseteq \C.
    \end{align}
    This is termed the \eax{field of algebraic numbers}.
\end{definition}
Of course, the fact that $\overline{\Q}$ is a field is not immediate. We prove minor results to support this claim.

\begin{lemma}
    Let $\alpha \in \overline{\Q}$ and $m_{\alpha}(X)$ be the (a) non-zero minimal degree polynomial over $\Q$ with $m_{\alpha}(\alpha) = 0$. Then $m_{\alpha}(X)$ is irreducible over $\Q$.
\end{lemma}

\begin{lemma}
    Let $\alpha \in \overline{\Q}$. Then $\Q[\alpha] = \Q(\alpha)$, and $[\Q(\alpha):\Q]$ is finite.
\end{lemma}

The concept of algebraic numbers can be generalized to any field $k$. Let $k \subseteq L$ be a field extension. An element $\alpha \in L$ is said to be algebraic over $k$ if $\alpha$ is a root of a non-zero polynomial over $k$. Here, $k[\alpha] = k(\alpha)$, and $[k(\alpha):k] = \deg m_{\alpha(X)}$ is finite.

\begin{lemma}
    Let $E \supseteq L \supseteq k$ be fields. Then $[E:k] = [E:L][L:k]$.
\end{lemma}

\begin{proof}
    If one of $E/L$ or $L/k$ is infinite dimensional, then $E/k$ is infinite dimensional and the equality holds. Now assume $E/k$ is finite dimensional. Clearly, $[E:k]$ being of finite dimension implies that $[E:L]$ and $[L:k]$ are finite dimensional. Let $\{e_{1},\ldots,e_{m}\}$ be an $L$-bases for $E$, and $\{\ell_{1},\ldots,\ell_{n}\}$ be a $k$-basis for $L$. We claim that $\{e_{i}\ell_{j} \mid 1 \leq i \leq m, 1 \leq j \leq n\}$ is a $k$-basis for $E$. It is clear that this set spans $E$ over $k$. To show linear independence, assume $\sum_{i,j} a_{ij} e_{i}\ell_{j} = 0$ for some $a_{ij} \in k$. Then $\sum_{i} (\sum_{j} a_{ij}\ell_{j}) e_{i} = 0$. Since $\{e_{1},\ldots,e_{m}\}$ is linearly independent over $L$, we have $\sum_{j} a_{ij}\ell_{j} = 0$ for each $i$. Since $\{\ell_{1},\ldots,\ell_{n}\}$ is linearly independent over $k$, we have $a_{ij} = 0$ for all $i,j$. Thus, the set is linearly independent over $k$, and hence a basis. Therefore, $[E:k] = mn = [E:L][L:k]$.
\end{proof}

\begin{definition}
    A field extension $L/k$ is called an \eax{algebraic extension} if all elements of $L$ are algebraic over $k$. 
\end{definition}

\begin{theorem}
    Every finite extension of a field is algebraic.
\end{theorem}
\begin{proof}
    Let $L/k$ be a finite extension. Then for any $\alpha \in L$, set the left multiplication map $\ell_{\alpha} : L \to L$ defined by $\ell_{\alpha}(x) = \alpha x$. This is a $k$-linear map, and hence has a characteristic polynomial $P_{\alpha}(X)$. By the Cayley-Hamilton theorem, $P_{\alpha}(\ell_{\alpha}) = 0$, and hence $P_{\alpha}(\alpha) = 0$. Thus, $\alpha$ is algebraic over $k$.
\end{proof}

\textit{July 22nd.}

\begin{theorem}
    Let $E \supseteq L \supseteq k$ be fields. Then $E/k$ is algebraic if and only if both $E/L$ and $L/k$ are algebraic.
\end{theorem}

Such an inclusion is termed a \eax{tower of extensions}. To prove this theorem, we require the following lemma.

\begin{lemma}
    Let $L \supseteq k$ be a field extesion, and let $\alpha_{1},\ldots,\alpha_{n} \in L$ be algebraic over $k$. Then $k(\alpha_{1},\ldots,\alpha_{n})$ is finite dimensional, and algebraic, over $k$.
\end{lemma}

A field extension $L = k(\beta)$ is called a \eax{simple extension}. In general, if $L/k$ is a field extension, then $L/k$ is called a \eax{finitely generated extension} if there exists a finite subset $S \subseteq L$ such that $L = k(S)$. Note that a finitely generated extension need not be a finite extension. For example, $\Q(\pi)/\Q$ is finitely generated but not finite.

\begin{proof}[Proof of theorem]
    Let us assume $L/E$ and $E/k$ are algebraic; let $\alpha \in E$. Then $\alpha$ is algebraic over $L$, and there exists a non-zero polynomial $g(X) \in L[X]$ such that $g(\alpha) = 0$. Write $g(X) = a_{0} + a_{1}x + \cdots + a_{r}X^{r}$ with $a_{i} \in L$. Now consider $k(a_{0},a_{1},\ldots,a_{r})$; here, each $a_{i}$ is algebraic over $k$ so $k(a_{0},a_{1},\ldots,a_{r})/k$ is algebraic. Consider $k(a_{0},a_{1},\ldots,a_{r})(\alpha)$; since $\alpha$ is algebraic over $k(a_{0},a_{1},\ldots,a_{r})$, we have $[k(a_{0},a_{1},\ldots,a_{r})(\alpha) : k(a_{0},a_{1},\ldots,a_{r})] < \infty$. Thus, $[k(a_{0},a_{1},\ldots,a_{r})(\alpha) : k]$ is finite, and hence $\alpha$ is algebraic over $k$. Since $\alpha \in E$ was arbitrary, we have that $E/k$ is algebraic. The converse is immediate since any $\alpha \in L$ is also in $E$, and hence algebraic over $k$.
\end{proof}

\begin{example}
    As an exercise, show that an algebraic extension $L/k$ is finitely generated if and only if $[L:k]$ is finite.
\end{example}

\section{Number Fields}

Simply put, a finite field extension of the rationals $\Q$ is called a \eax{number field}. Examples include $\Q(\sqrt{2})$, $\Q(\alpha_{1},\ldots,\alpha_{r})$ where $\alpha_{i}$ are algebraic over $\Q$. Non-examples include $\Q(\pi)$, $\Q(\sqrt{2},\sqrt{3},\sqrt{5},\ldots,\sqrt{p},\ldots)$ for all primes $p$. Recall the field of algebraic numbers:
\begin{align}
    \overline{\Q} = \{\alpha \in \C \mid \alpha \text{ is algebraic over } \Q\}.
\end{align}
It is left as an exercise to show that $\overline{\Q}$ is a field, and that $\overline{\Q}/\Q$ is an algebraic extension.

\begin{definition}
    A field $K$ is called \eax{algebraically closed} if every non-constant polynomial $f(X) \in K[X]$ has a root in $K$. Equivalently, irreducible polynomials over $K$ are exactly the linear polynomials.
\end{definition}
For example, $\C$ is algebraically closed, but $\R$ is not. This is a consequence of the fundamental theorem of algebra.

\begin{example}
    The field $\overline{\Q}$ is algebraically closed. To show this, let $f(X) \in \overline{\Q}[X]$ be an irreducible polynomial. By the fundamental theorem of algebra, there exists a root $\alpha \in \C$ of $f(X)$. So, $\overline{\Q}(\alpha)/\overline{\Q}$ is algebraic. But since $\overline{\Q}/\Q$ is algebraic, we have that $\overline{\Q}(\alpha)/\Q$ is algebraic. Thus, $\alpha \in \overline{\Q}$, and hence $f(X)$ has a root in $\overline{\Q}$.
\end{example}

If we define $\overline{\Q}^{\R} = \{\alpha \in \R \mid \alpha \text{ is algebraic over } \Q\}$. Clearly, $\overline{\Q}^{\R}$ is algebraic. Consider $i \in \C$; $i$ is algebraic over $\Q$ but $i \notin \R$, so $i \notin \overline{\Q}^{\R}$. Thus, the minimal polynomial of $i$ over $\overline{\Q}^{\R}$ is $X^{2}+1$, which is irreducible over $\overline{\Q}^{\R}$. Hence, $\overline{\Q}^{\R}$ is not algebraically closed. 
\\

\noindent \textit{July 23rd.}

Let $\varphi:F \to R$ be a ring homomorphism between a non-trivial field $F$ and a non-trivial ring $R$. Since $F$ is a field, it has only two ideals, $\{0\}$ and $F$. Enforcing that $\varphi(1) \neq 0$, we find that $\ker \varphi = \{0\}$, and hence $\varphi$ is injective. Hence, on various occasions, we will identify $F$ with its image $\varphi(F) \subseteq R$. 

\begin{theorem}
    Let $k$ be a field. Then there exists an algebraically closed field $K$ containing $k$.
\end{theorem}
To provide a sketch of a proof, pick, for each non-constant polynomial $f(X) \in k[X]$, a variable $X_{f}$. Consider the polynomial ring $R = k[X_{f} \mid f \in k[X], f \text{ non-constant}]$. Let $I$ be the ideal generated by $\{f(X_{f}) \mid f \in k[X], f \text{ non-constant}\}$. This $I$ is a proper ideal since it does not contain the unity. By Zorn's lemma, there exists a maximal ideal $\mf{m} \supseteq I$. Now define $K = R/\mf{m}$, a field. Since $k \hookrightarrow R \to K$, we have that $K$ is a field extension of $k$. This $K$ is \textit{not} equal to $k$ since the element $X_{f} + \mf{m}$ is not in $k$ but is algebraic over $k$ since $f(X_{f} + \mf{m}) = 0$ in $K_{0}$. Hence, $K/k$ is an algebraic extension. It now follows that every $f(X) \in k[X]$ with degree at least 1 has a root in $K = K_{0}$. Repeating this process, we can construct a tower of algebraic extensions $K_{0} \subseteq K_{1} \subseteq K_{2} \subseteq \cdots$, and define $E = \bigcup_{n=0}^{\infty} K_{n}$. Then $E/k$ is algebraic, and every non-constant polynomial in $E[X]$ has a root in $E$. Hence, $E$ is algebraically closed.

\begin{definition}
    Let $k$ be a field. Then the \eax{algebraic closure} of $k$ is an algebraically closed field $\overline{k}$ containing $k$ such that $\overline{k}/k$ is algebraic.
\end{definition}

We know that if $L \supseteq k$ is a field extension, then $L$ is also a $k$-vector space. For any algebraic closure $\overline{k}$ of $k$ (or an algebraically closed field containing $k$), a homomorphism $\varphi:L \to \overline{k}$ is called a \eax{$k$-embedding} if $\varphi$ is $k$-linear homomorphism of fields.

\begin{example}
    We look at some examples of $k$-embeddings.
    \begin{itemize}
        \item Let $k = \R$ and $\overline{k} = \C$. With $L = \R(\sqrt{-2})$, which is essentially $\C$, calling $\alpha = \sqrt{-2}$, we have that $\varphi(a+b\alpha) = a+b\overline{\alpha}$ and $\varphi(a+b\alpha) = a+b\alpha$ are the only two $\R$-embeddings of $L$ into $\C$.
        \item Let $\alpha$ be algebraic over $k$, and let $E = k(\alpha)$. We ask what are all the $k$-embeddings of $E$ into $\overline{k}$. Let $m_{\alpha}(X)$ be the minimal polynomial of $\alpha$ over $k$. Pick any root $\beta \in \overline{k}$ of $m_{\alpha}(X)$, and let $\deg m_{\alpha}(X) = d$. Define a mapping $\varphi:k(\alpha) \to \overline{k}$ by $\alpha \mapsto \beta$, and extending $k$-linearly. Then $\varphi$ is a $k$-embedding of $E$ into $\overline{k}$. Conversely, if $\sigma:k(\alpha) \to \overline{k}$ is a $k$-embedding, then $\sigma(\alpha)$ is a root of $m_{\alpha}(X)$ in $\overline{k}$. Thus, the $k$-embeddings of $E$ into $\overline{k}$ are in bijection with the roots of $m_{\alpha}(X)$ in $\overline{k}$.
    \end{itemize}
\end{example}